\section{Limitations}
\label{sec:limitations}

The most immediate limitation is the absence of pretrained-model results.
The current implementation and protocol cannot support a claim of improved
hallucination, generalization, efficiency, or competitive performance.
Comparison to close methods and verification of adapted baselines remain
substantial experimental requirements.

Architecture support does not establish baseline parity. Our GrAInS variant
holds Qwen3 DeepStack features and Gemma per-layer inputs fixed during
first-stream attribution. The current VTI port also differs from author code
in normalization, masking, and PCA. These departures require validation before
any claim of superiority over the published methods.

The localization intervention is also imperfect. Gray occlusion creates a
distribution shift, and a translated background mask cannot match every
property of an object-region edit. Unannotated objects and small residual
target fragments can invalidate labels. Semantic-context edits introduce
their own artifacts, so conclusions must also hold on natural images.
Even a successful controlled comparison does not uniquely attribute the
failure to a language prior rather than a visual co-occurrence shortcut.

Selected heads are aggregates of many computations. A head-level causal
effect does not imply a monosemantic feature or a complete circuit.
Moreover, selection is based on positive clean examples and margin drops;
the resulting sample may underrepresent difficult targets. Fixed gains
can act differently across questions and image conditions. A useful oracle
patching effect does not establish that a fixed deployable gain exists.

Finally, preserving the next-token distribution on correct development
examples does not guarantee preservation of multi-token descriptions,
counting, spatial relations, or reasoning. Our claims are restricted to the
tested tasks and controlled data. The study uses public image datasets
under their applicable terms, and neither diagnostic edits nor aggregate
metrics should be presented as a general reliability guarantee.
